% ============================================================
%  PLANTILLA — Reporte Etapa 2
%  Estructura mapeada 1:1 a la rúbrica del curso (45%).
%
%  Cada sección lleva en comentario el % que pesa y los
%  criterios "Excelente" que tiene que cubrir.
%  Reemplazar los TODOs y comentarios con contenido real.
% ============================================================

\documentclass[11pt,a4paper]{article}

% ---------- Paquetes ----------
\usepackage[utf8]{inputenc}
\usepackage[T1]{fontenc}
\usepackage[spanish]{babel}
\usepackage[margin=2.5cm]{geometry}
\usepackage{amsmath,amssymb}
\usepackage{booktabs}
\usepackage{graphicx}
\usepackage{subcaption}
\usepackage{hyperref}
\usepackage{xcolor}
\usepackage{tcolorbox}
\usepackage{fancyhdr}
\usepackage{caption}
\usepackage{float}
\usepackage{listings}

\hypersetup{
  colorlinks=true,
  linkcolor=blue!50!black,
  citecolor=blue!50!black,
  urlcolor=blue!50!black,
}

\lstset{
  basicstyle=\ttfamily\small,
  frame=single,
  breaklines=true,
  showstringspaces=false,
  extendedchars=true,
  inputencoding=utf8,
  literate=%
    {á}{{\'a}}1 {é}{{\'e}}1 {í}{{\'i}}1 {ó}{{\'o}}1 {ú}{{\'u}}1
    {Á}{{\'A}}1 {É}{{\'E}}1 {Í}{{\'I}}1 {Ó}{{\'O}}1 {Ú}{{\'U}}1
    {ñ}{{\~n}}1 {Ñ}{{\~N}}1 {ü}{{\"u}}1 {Ü}{{\"U}}1
    {¿}{{?`}}1 {¡}{{!`}}1,
}

% --- Placeholder para figuras ausentes (Overleaf sin assets locales) ---
% Uso: \safeincludegraphics[opciones]{ruta}
% Si el archivo existe -> \includegraphics normal.
% Si no -> caja gris con el path para que el documento compile igual.
\newcommand{\safeincludegraphics}[2][width=0.7\textwidth]{%
  \IfFileExists{#2}{%
    \includegraphics[#1]{#2}%
  }{%
    \fbox{\parbox[c][4cm][c]{0.7\textwidth}{\centering\ttfamily\small
      [imagen pendiente]\\\detokenize{#2}}}%
  }%
}

\begin{document}

% --- Bloque de título estilo REVTeX/AIP ---
\begin{center}
  {\LARGE\bfseries Modelos de Espacio de Estados (Mamba)\\
   para \textit{Vetting} de Exoplanetas en TESS\par}
  \vspace{1.2em}
  {\normalsize José F. Zumbado Ruiz and Jeremmy Aguilar Villanueva\par}
  \vspace{0.4em}
  {\itshape Escuela de Ingeniería en Computación, Instituto Tecnológico
   de Costa Rica, Cartago, Costa Rica\par}
  \vspace{0.3em}
  {(Reporte Etapa 2, IC8104 Inteligencia Artificial, Semestre I 2026)\par}
  \vspace{0.1em}
  {(Dated: \today)\par}
\end{center}
\vspace{1.5em}

\begin{abstract}
% TODO (150-200 palabras). Estructura sugerida:
% (1) Problema: vetting de TOIs en TESS, dataset 1576 ejemplos.
% (2) Aporte: comparación controlada de 4 modelos (Random, LogReg, CNN, Mamba).
% (3) Resultado principal: Mamba single AUC test = X ± Y, supera CNN por Z pp.
% (4) Multi-seed sweep + ensemble reportado como métrica final.
% (5) Reproducibilidad: configs YAML + seeds + test sellado.
\end{abstract}

\tableofcontents
\newpage

% ============================================================
\section{Introducción}
\label{sec:intro}

\subsection{Motivación}
\label{subsec:motivacion}

La misión TESS (\textit{Transiting Exoplanet Survey Satellite}) de la NASA
fue diseñada para observar a cadencia de 2 minutos las $\sim$200\,000
estrellas enanas más brillantes del cielo, además de tomar imágenes de
campo completo (FFI) que generan curvas adicionales para muchas otras
estrellas. Cada sector observado dura $\sim$27 días y produce series
temporales del orden de $L \approx 18\,000$ puntos. De ese flujo de
fotometría, una pequeña fracción muestra señales candidatas a tránsito
planetario; la NASA cataloga estos candidatos como TOIs
(\textit{TESS Objects of Interest}).

El catálogo TOI consultado para este trabajo registra más de 7\,000
candidatos según el \textit{snapshot} utilizado (al 1 de julio de 2025
se reportaban 7\,655 TOIs, 638 con etiqueta CP y 1\,987 con etiqueta FP).
Por el costo de la revisión humana, el \textit{vetting} de TOIs se apoya
cada vez más en clasificadores automáticos.

El \textit{dataset} usado en este proyecto se restringe a TOIs con
etiqueta CP (\textit{Confirmed Planet}) o FP (\textit{False Positive}).
Se excluye KP (\textit{Known Planet}, planeta confirmado por misiones
previas) por decisión experimental, para evitar mezclar planetas
conocidos pre-TESS con confirmaciones obtenidas por la propia misión;
se excluyen también PC (\textit{Planet Candidate}) y APC (\textit{Ambiguous Planet Candidate}) por carecer de etiqueta definitiva. Tras
filtrar TICs sin cobertura fotométrica descargable mediante
\texttt{lightkurve}, se obtiene el conjunto Tier 1 de 1\,576 TICs
etiquetados. 

Para situar este tamaño respecto a los baselines de la literatura,
AstroNet \cite{shallue2018} se entrenó sobre 15\,737 TCEs etiquetados
de Kepler (3\,600 PC, 9\,596 AFP y 2\,541 NTP), y ExoMiner
\cite{exominer} reporta un \textit{dataset} de trabajo de 30\,609
ejemplos sobre $\sim$34\,000 TCEs de Kepler DR25. El \textit{dataset}
de este proyecto es aproximadamente un orden de magnitud menor que el
de AstroNet y cerca de 20 veces menor que el de ExoMiner. Esta
diferencia de escala condiciona la magnitud absoluta de las métricas
reportadas y se discute al evaluar la comparación con la literatura
(Sec.\ref{sec:discussion}).

Los modelos de referencia en la literatura, como AstroNet
\cite{shallue2018} y ExoMiner \cite{exominer}, utilizan arquitecturas
profundas basadas en representaciones de curvas de luz, con componentes
convolucionales y vistas globales y locales. Este enfoque tiene dos
limitaciones documentadas. Primero, el campo receptivo de una CNN crece
de forma lineal con la profundidad y la captura de dependencias largas
en secuencias de $L=18\,000$ requiere arquitecturas profundas, costosas
de entrenar. Segundo, las arquitecturas tipo Transformer, que sí modelan
dependencias largas de forma nativa, tienen complejidad $\mathcal{O}(L^2)$
en memoria y cómputo, lo cual es prohibitivo bajo la restricción de
hardware del proyecto (RTX 3050 con 4\,GB de VRAM).

Los modelos de espacio de estados selectivos (Selective State Space Models,
SSM) propuestos por Gu y Dao \cite{gu2023mamba} ofrecen complejidad
$\mathcal{O}(L)$ con memoria constante mediante un \textit{selective scan},
y han igualado o superado a los Transformers en varias tareas de secuencia
larga. Hasta donde se revisó en la literatura del proyecto, no se
encontró una aplicación previa de Mamba al \textit{vetting} de TOIs en
TESS bajo restricciones de hardware similares. Esta brecha constituye
la motivación técnica del proyecto.

\subsection{Objetivos de la Etapa 2}
\label{subsec:objetivos}

La Etapa 2 del curso evalúa la fase de
\textit{Modelado, Entrenamiento, XAI y Evaluación}. Los criterios de la
rúbrica fijan cinco ejes de trabajo. Este reporte se organiza para
cubrir cada uno de forma explícita.

\begin{enumerate}
  \item \textbf{Baselines obligatorios.} Entregar al menos dos baselines
    razonables y una comparación clara contra el modelo principal. En este
    trabajo se construye una \textit{escalera} de cuatro modelos
    (Sec.\ref{sec:baselines}) que permite atribuir cada incremento de
    desempeño a un componente concreto de la arquitectura.
  \item \textbf{Validación y protocolo.} Adoptar un protocolo sólido con
    \textit{test} ciego, \textit{split} por TIC y prevención explícita
    de \textit{leakage} (Sec.\ref{sec:protocol}).
  \item \textbf{Optimización de hiperparámetros.} Realizar una búsqueda
    sistemática con justificación de costo (Sec.\ref{sec:tuning}). Por
    restricciones de VRAM se sustituye la validación cruzada por una
    estrategia multi-semilla equivalente en términos de variabilidad
    estadística. Cada fold adicional replica íntegramente el entrenamiento, lo que con K=5 y un
   sweep de ~20 configuraciones escalaría a más de 100 horas de GPU en una RTX
  3050 de 4 GB VRAM, haciendo la validación cruzada computacionalmente inviable 
  en el contexto de este proyecto.
  \item \textbf{Métricas y análisis de resultados.} Reportar métricas
    apropiadas a la clase desbalanceada, curvas ROC y PR, análisis de
    errores y variabilidad seed-by-seed (Secs.~\ref{sec:results}
    y~\ref{sec:errors}).
  \item \textbf{Interpretabilidad (XAI).} Aplicar técnicas adecuadas al
    tipo de modelo y derivar hallazgos accionables (Sec.\ref{sec:xai}).
    Dado que Mamba no implementa atención en el sentido del Transformer,
    se utilizan métodos basados en gradiente y oclusión en lugar de
    mapas de atención.
\end{enumerate}

A esto se suma el eje transversal de \textbf{ingeniería y
reproducibilidad} (Sec.\ref{sec:repro}), cubierto mediante configuraciones
YAML por experimento, semillas fijas, un protocolo de \textit{test} sellado
y una suite de pruebas automatizadas (Sec.\ref{sec:tests}).

\subsection{Contribuciones}
\label{subsec:contribuciones}

Este reporte aporta lo siguiente.

\begin{enumerate}
  \item Una \textbf{comparación controlada de cuatro modelos} sobre el mismo
    \textit{split} sellado de TESS: \textit{random} estratificado, regresión
    logística sobre el catálogo TOI, CNN \textit{single-branch} inspirada
    en AstroNet \cite{shallue2018}, y Mamba con vista global como modelo
    principal.
  \item Una \textbf{exploración experimental de Mamba aplicado al
    \textit{vetting} de TOIs en TESS}. Hasta donde se revisó en la
    literatura del proyecto, no se encontró una aplicación previa bajo
    restricciones de hardware similares. El modelo se entrena en FP16
    con \textit{gradient checkpointing}, lo que permite procesar
    secuencias de longitud 18\,000 dentro de 4\,GB de VRAM.
  \item Un \textbf{protocolo multi-semilla} (cinco semillas para Mamba
    \textit{single-view}) que reemplaza la validación cruzada $K$-fold
    sin sacrificar la estimación de variabilidad. La sustitución se
    justifica en términos del costo estimado y se documenta en
    Sec.\ref{sec:tuning}.
  \item Un \textbf{análisis XAI multi-método} (\textit{gradient saliency},
    \textit{integrated gradients} y \textit{occlusion sensitivity}) sobre
    Mamba, con una política de selección reproducible de ocho casos
    (los dos más confidentes por cuadrante TP, TN, FN, FP).
  \item Una \textbf{infraestructura de reproducibilidad} que incluye
    configuraciones YAML versionadas, semillas controladas, una suite de
    49 pruebas automatizadas, manifiestos de entorno por corrida y un
    protocolo de \textit{test} sellado evaluado una sola vez por modelo.
\end{enumerate}

Como insumo para Etapa 3 del curso (\textit{paper} IEEE/ACM final), el
repositorio incluye además modelos exploratorios adicionales cuyos
resultados están documentados en \texttt{paper/results/} pero quedan
fuera del alcance del presente reporte.



% ============================================================
\section{Datos y preprocesamiento}
\label{sec:data}

\subsection{Fuentes y construcción del dataset}
\label{subsec:fuentes}

El \textit{dataset} se construye a partir de dos fuentes públicas de la NASA
que se consultan programáticamente.

\begin{itemize}
  \item \textbf{NASA Exoplanet Archive: TOI Catalog.} Tabla pública con
    los TOIs (\textit{TESS Objects of Interest}) y su disposición
    científica (CP, FP, PC, KP, APC, FA). Se descarga vía una consulta
    TAP (\textit{Table Access Protocol}) ejecutada por
    \texttt{scripts/get\_data.py}. Del catálogo se conservan únicamente
    las columnas relevantes para el experimento: identificador
    (\texttt{tid}), disposición (\texttt{tfopwg\_disp}), y los parámetros
    de tránsito utilizados como \textit{features} por la regresión
    logística (Sec.\ref{subsec:logreg}): \texttt{pl\_orbper},
    \texttt{pl\_trandurh}, \texttt{pl\_trandep} y \texttt{st\_tmag}.
  \item \textbf{MAST Archive.} Almacén oficial de los archivos FITS
    producidos por la \textit{pipeline} SPOC (\textit{Science Processing
    Operations Center}) de TESS. Para cada TIC se descarga la curva
    \texttt{PDCSAP\_FLUX} a cadencia de 2 minutos, que ya viene
    detrendeada por la \textit{pipeline} oficial. La descarga se realiza
    con la librería \texttt{lightkurve}, mediante el \textit{script}
    \texttt{scripts/download\_lightcurves.py}, idempotente y con
    \textit{manifest} versionado (\texttt{data/splits/manifest.csv}).
\end{itemize}

La construcción del \textit{dataset} sigue una secuencia de filtros
declarada explícitamente:

\begin{enumerate}
  \item Del catálogo TOI se conservan únicamente las disposiciones
    CP (\textit{Confirmed Planet}, etiqueta positiva) y FP
    (\textit{False Positive}, etiqueta negativa). Se excluye KP por
    decisión experimental (Sec.\ref{subsec:motivacion}); se excluyen
    PC y APC por carecer de etiqueta definitiva.
  \item Se descartan TICs sin \texttt{PDCSAP\_FLUX} accesible en MAST
    o cuya descarga vía \texttt{lightkurve} falla tras los reintentos
    configurados.
  \item Se descartan TICs cuyo preprocesamiento global produce un
    tensor degenerado (por ejemplo, curva con más del 50\,\% de
    \texttt{NaN} o longitud útil menor a 5\,000 puntos).
\end{enumerate}

El conjunto resultante contiene 1\,576 TICs etiquetados. Sus particiones
\textit{train}, \textit{validation} y \textit{test} se construyen por TIC
(Sec.\ref{sec:protocol}) en proporción aproximada 70/15/15.
Tabla~\ref{tab:dataset-counts} resume los conteos.

\begin{table}[H]
\centering
\caption{Conteos por \textit{split} y clase. El \textit{split} se
realiza por TIC y de forma estratificada para preservar la proporción
CP/FP entre particiones (Sec.\ref{sec:protocol}).}
\label{tab:dataset-counts}
\begin{tabular}{lccc}
\toprule
\textit{Split} & N & CP & FP \\
\midrule
Train      & 1\,103 & 422 & 681 \\
Validation & 236    & 90  & 146 \\
Test       & 237    & 91  & 146 \\
\midrule
Total      & 1\,576 & 603 & 973 \\
\bottomrule
\end{tabular}
\end{table}

\subsection{Vista global}
\label{subsec:global-view}

La \textbf{vista global} es la representación principal del modelo
Mamba y de las CNN \textit{single-branch}. Se construye en
\texttt{scripts/preprocess\_global.py} con tres pasos:

\paragraph{1. Lectura de la señal cruda.}
Por cada TIC se concatenan los flujos \texttt{PDCSAP\_FLUX} de todos
los sectores descargados, ordenados por tiempo. Esta señal ya viene
detrendeada por SPOC y con artefactos sistemáticos removidos.

\paragraph{2. Normalización por curva individual.}
Sea $\mathbf{f} = (f_1, \dots, f_T)$ la señal cruda concatenada y
$\tilde{f} = \mathrm{median}(\mathbf{f})$ su mediana. La curva
normalizada $\mathbf{x}$ se define como
\begin{equation}
  x_t = \frac{f_t}{\tilde{f}}, \qquad t = 1, \dots, T.
\label{eq:median-norm}
\end{equation}
La normalización se hace siempre con estadísticas de la propia curva,
nunca con estadísticas globales del \textit{dataset}. Esta decisión
evita el \textit{leakage} de escala desde el \textit{train} al
\textit{test} (Sec.\ref{sec:protocol}).

\paragraph{3. Ajuste a longitud fija $L=18\,000$.}
Las curvas resultantes tienen longitudes variables (función del número
de sectores observados). Para que el modelo opere con tensores de forma
homogénea se aplica:
\begin{itemize}
  \item Si $T \geq L$, se recorta al primer bloque contiguo de $L$
    puntos válidos.
  \item Si $T < L$, se rellena con la mediana de la curva (es decir,
    con el valor neutro $1.0$ tras la normalización).
\end{itemize}
Los valores \texttt{NaN} restantes se interpolan linealmente sobre el
índice temporal. La salida es un tensor \texttt{float32} de forma
$(1, 18\,000)$ guardado como
\texttt{data/processed/global/<tic>.pt}. Este contrato lo verifica
\texttt{test\_dataset.py} (Sec.\ref{sec:tests}).

\paragraph{Qué modelos la consumen.}
La vista global es el \textit{input} de los modelos de aprendizaje
profundo evaluados en este reporte: la CNN \textit{single-branch}
(Sec.\ref{subsec:cnn}) y Mamba \textit{single-view}
(Sec.\ref{subsec:mamba}). No se construyen vistas adicionales sobre la
curva.

\subsection{Data augmentation}
\label{subsec:augmentation}

Para mitigar el tamaño reducido del \textit{dataset}
(Sec.\ref{subsec:motivacion}) se aplican cuatro transformaciones
estocásticas a la vista global, definidas en
\texttt{src/exoplanet/data/augment.py}. Estas se aplican
\textbf{únicamente al \textit{split} de \textit{train}}; \textit{val}
y \textit{test} se evalúan con la señal cruda preprocesada, para que
las métricas reflejen desempeño sobre datos no modificados.

\begin{table}[H]
\centering
\caption{Transformaciones de \textit{augmentation} aplicadas al
\textit{split} de \textit{train}.}
\label{tab:augmentation}
\renewcommand{\arraystretch}{1.25}
\begin{tabular}{p{0.22\textwidth} p{0.32\textwidth} p{0.38\textwidth}}
\toprule
\textbf{Transformación} & \textbf{Operación} & \textbf{Hiperparámetros del proyecto} \\
\midrule
\texttt{temporal\_shift} &
  Desplazamiento circular de la curva: $x'_t = x_{(t + s) \bmod L}$
  con $s \sim \mathcal{U}\{-S, \dots, S\}$. &
  $S = 500$ (puntos). \\
\addlinespace
\texttt{gaussian\_noise} &
  $x'_t = x_t + \varepsilon_t,\ \varepsilon_t \sim \mathcal{N}(0, \sigma^2)$. &
  $\sigma = 0.001$ (orden del ruido instrumental TESS típico). \\
\addlinespace
\texttt{time\_reverse} &
  Inversión temporal $x'_t = x_{L-t+1}$ con probabilidad $p$.
  Válida por la simetría temporal del tránsito. &
  $p = 0.5$. \\
\addlinespace
\texttt{amplitude\_scale} &
  $x'_t = a \cdot x_t,\ a \sim \mathcal{U}(a_{\min}, a_{\max})$.
  Simula diferencias de calibración. &
  $a \in [0.99, 1.01]$. \\
\bottomrule
\end{tabular}
\end{table}

Las cuatro transformaciones cumplen el mismo contrato funcional:
son puras (no mutan el tensor de entrada), preservan forma y
\texttt{dtype}, son determinísticas bajo \texttt{torch.Generator}, y
manejan correctamente los casos identidad (\textit{e.g.}, $\sigma=0$
no agrega ruido y devuelve un clon del tensor). La validación de este
contrato se hace mediante 21 pruebas unitarias en
\texttt{tests/test\_augment.py} (Sec.\ref{sec:tests}). El orden de
aplicación y los hiperparámetros se declaran por experimento en el
campo \texttt{augmentation} del YAML correspondiente.

% ============================================================
\section{Baselines}
\label{sec:baselines}

Para evaluar el aporte real de la arquitectura Mamba en este problema se
construye una \textbf{escalera de cuatro modelos} ordenada por
expresividad creciente. La idea detrás de la escalera es que cada
modelo agrega exactamente un componente respecto al anterior, de modo
que la diferencia de desempeño entre dos vecinos puede atribuirse a
ese componente concreto y no a una combinación de cambios. La
Tabla~\ref{tab:baselines-roles} resume los cuatro modelos y su rol.

\begin{table}[H]
\centering
\caption{Escalera de modelos. La columna \textit{Aporta} describe el
componente que cada modelo agrega respecto al anterior; el orden
permite atribuir ganancias o pérdidas a una variable a la vez.}
\label{tab:baselines-roles}
\renewcommand{\arraystretch}{1.2}
\begin{tabular}{p{0.27\textwidth} p{0.30\textwidth} p{0.36\textwidth}}
\toprule
\textbf{Modelo} & \textbf{Aporta respecto al anterior} & \textbf{Rol en el reporte} \\
\midrule
Random estratificado &
  Solo el \textit{prior} de clase &
  Piso mínimo (Sec.\ref{subsec:random}) \\
\addlinespace
Regresión logística (catálogo) &
  Features tabulares del TOI &
  Baseline no-DL fuerte (Sec.\ref{subsec:logreg}) \\
\addlinespace
CNN single-branch &
  Aprende representación 1D de la vista global &
  Baseline DL clásico (Sec.\ref{subsec:cnn}) \\
\addlinespace
\textbf{Mamba single-view} &
  Sustituye convolución por SSM selectivo &
  \textbf{Modelo principal de Etapa 2} (Sec.\ref{subsec:mamba}) \\
\bottomrule
\end{tabular}
\end{table}

El modelo entregable principal de Etapa 2 es Mamba \textit{single-view}.
Los tres modelos anteriores establecen una jerarquía creciente de
expresividad contra la cual se mide su aporte real.

\subsection{Random estratificado}
\label{subsec:random}

El baseline más simple. No mira la curva: solo conoce la proporción
de clases en el \textit{split} de entrenamiento. Para cada muestra de
test predice la clase positiva con probabilidad
\begin{equation}
\hat{p}(y=1) = \frac{N_{\text{CP}}^{\text{train}}}{N_{\text{CP}}^{\text{train}} + N_{\text{FP}}^{\text{train}}} = \frac{422}{1\,103} \approx 0.383,
\label{eq:random-prior}
\end{equation}
y emite una predicción de Bernoulli con ese parámetro.

\paragraph{Para qué sirve.}
Establece el \textit{piso} contra el cual todos los modelos posteriores
deben superar. Por construcción su AUC-ROC esperado es exactamente
$0.5$, ya que clasifica sin información alguna sobre la curva. Si un
modelo no supera este piso, es evidencia de un error de implementación
o de pipeline, no de un problema de aprendizaje.

\subsection{Regresión logística sobre catálogo}
\label{subsec:logreg}

El segundo baseline aprovecha la información tabular ya presente en el
catálogo TOI sin tocar las curvas de luz. Por cada TIC se construye un
vector de \textit{features}
\(\mathbf{x} = (\texttt{pl\_orbper},\,\texttt{pl\_trandurh},\,\texttt{pl\_trandep},\,\texttt{st\_tmag}, \dots)\),
y se ajusta un modelo de regresión logística:
\begin{equation}
\hat{p}(y=1 \mid \mathbf{x}) = \sigma(\mathbf{w}^\top \mathbf{x} + b),
\qquad \sigma(z) = \frac{1}{1 + e^{-z}}.
\label{eq:logreg-sigmoid}
\end{equation}
La función $\sigma$ (sigmoide) comprime cualquier valor real a un
número entre 0 y 1, interpretable como probabilidad. Los parámetros
$\mathbf{w}, b$ se aprenden minimizando la pérdida de entropía cruzada
binaria con pesos por clase para corregir el desbalance:
\begin{equation}
\mathcal{L}_{\text{BCE}}(\mathbf{w}, b) = -\frac{1}{N}\sum_{i=1}^{N}
  \Bigl[ w_+\,y_i\log \hat{p}_i + (1-y_i)\log(1-\hat{p}_i) \Bigr],
\quad w_+ = \frac{N_{\text{FP}}}{N_{\text{CP}}}.
\label{eq:bce-weighted}
\end{equation}

\paragraph{Decisión de diseño.}
Las \textit{features} usadas vienen \textbf{directamente del catálogo},
no se recalculan desde la curva (por ejemplo, no corremos un algoritmo
BLS propio para reestimar el período). Esta elección se documenta para
ser transparentes: representa un compromiso entre rigor y costo.
Recalcular BLS sobre 1\,576 curvas con grilla densa toma horas y no era
parte del alcance de Etapa 2; el catálogo TOI ya tiene esos números
con la calidad oficial de la NASA.

\paragraph{Para qué sirve.}
Establece cuánto desempeño se puede extraer de la metadata sin tocar
la curva. Si un modelo basado en la curva no supera a la regresión
logística, significa que la curva no aporta señal adicional sobre lo
que ya está en el catálogo, lo cual sería un hallazgo paper-worthy en
sí mismo.

\subsection{CNN single-branch (AstroNet-inspired)}
\label{subsec:cnn}

Primera arquitectura de \textit{deep learning} de la escalera. Se
inspira en AstroNet (Shallue \& Vanderburg 2018) pero opera sobre una
sola rama (la vista global), lo cual la hace comparable
directamente con Mamba \textit{single-view}.

\paragraph{Operación básica de una capa convolucional 1D.}
Una capa convolucional aplica un filtro $\mathbf{k}$ de tamaño fijo
$K$ a la curva $\mathbf{x}$, deslizándolo a lo largo del tiempo:
\begin{equation}
y_t = \sum_{i=0}^{K-1} k_i \cdot x_{t-i} + b.
\label{eq:conv1d}
\end{equation}
El filtro tiene los mismos parámetros para todas las posiciones
temporales (\textit{weight sharing}), lo que le permite detectar el
mismo patrón local (por ejemplo, un \textit{dip} de tránsito) sin
importar dónde aparezca en la curva. Apilar varias capas
convolucionales con \textit{max-pooling} entre ellas agranda el
\textit{campo receptivo}: la profundidad $D$ de capas con \textit{kernel}
$K$ y \textit{pooling} $P$ produce un campo receptivo del orden de
$K \cdot P^D$ puntos.

\paragraph{Arquitectura usada.}
Cinco bloques convolucionales con canales crecientes
$(16, 32, 64, 128, 256)$, cada bloque con dos \texttt{Conv1d} + ReLU +
\texttt{MaxPool}, seguidos de \textit{flatten} y dos capas densas que
producen un logit. Total aproximado del orden de $10^5$ parámetros
entrenables. Esta variante se llama \textit{single-branch} porque solo
recibe \texttt{global\_view} como entrada; AstroNet original usa dos
ramas (global + local), simplificación que se documenta como
\textit{single-branch} en el código.

\paragraph{Limitación que motiva Mamba.}
Para secuencias de $L = 18\,000$ puntos, el campo receptivo de una CNN
poco profunda no abarca dependencias entre regiones distantes de la
curva (por ejemplo, dos tránsitos consecutivos separados por miles de
puntos). Capturarlas requiere o bien una CNN muy profunda (cara de
entrenar), o bien una arquitectura con dependencia explícitamente
secuencial: Transformers (de costo cuadrático) o SSM como Mamba.

\subsection{Mamba single-view}
\label{subsec:mamba}

\textbf{Modelo principal del reporte.} Mamba (Gu \& Dao 2023) pertenece
a la familia de \textit{Selective State Space Models} (SSM), una
generalización de las RNN clásicas con tres propiedades clave para
secuencias largas: (i) complejidad lineal $\mathcal{O}(L)$ en cómputo y
memoria, (ii) parámetros del estado que dependen del \textit{input}
(\textit{selective scan}), y (iii) implementación eficiente en GPU
mediante \textit{kernels} fusionados.

\paragraph{Intuición: estado oculto que recorre la secuencia.}
Imaginá un detective que va leyendo la curva de luz de izquierda a
derecha, anotando en un cuaderno
(el \textit{estado oculto} $h_t$) lo más relevante que ha visto.
En cada paso $t$ recibe el nuevo valor $x_t$ y decide qué de su
cuaderno actualizar y qué olvidar. Al final del recorrido emite un
juicio. Las RNN clásicas hacían esto pero olvidaban demasiado en
secuencias largas; Mamba mantiene un estado mucho más expresivo.

\paragraph{Ecuaciones del SSM selectivo.}
Sea $x_t \in \mathbb{R}$ el valor de la curva en el instante $t$ y
$h_t \in \mathbb{R}^N$ el estado oculto de dimensión $N$ (configurable
por el modelo). La dinámica continua es
$h'(t) = A\,h(t) + B\,x(t)$, $y(t) = C\,h(t)$, que se discretiza a
\begin{equation}
h_t = \bar{A}_t\,h_{t-1} + \bar{B}_t\,x_t,
\qquad
y_t = C_t\,h_t.
\label{eq:ssm}
\end{equation}
Lo \textbf{selectivo} de Mamba es que las matrices
$\bar{A}_t, \bar{B}_t, C_t$ \textit{dependen del input} $x_t$
mediante proyecciones lineales aprendidas: cada paso decide qué
información dejar pasar y qué descartar en función del propio valor
observado. Esto le da a Mamba una capacidad de filtrado dinámico que
las RNN clásicas no tenían.

\paragraph{Por qué importa para $L=18\,000$.}
Comparado con un Transformer, que tiene complejidad $\mathcal{O}(L^2)$
(cada token mira a todos los demás), Mamba tiene complejidad
$\mathcal{O}(L)$ porque el estado $h_t$ es de tamaño fijo y se
actualiza paso a paso. Para $L = 18\,000$:
\begin{itemize}
  \item Transformer sin trucos: $\sim L^2 = 3.24 \times 10^8$
    operaciones de \textit{attention} por capa. Memoria requerida
    para una matriz de atención FP16: $\sim 0.6\,\text{GB}$ por capa.
    \textbf{No cabe} en 4 GB de VRAM con $L = 18\,000$ a tamaños de
    \textit{batch} útiles.
  \item Mamba: $\sim L \cdot N = 1.4 \times 10^6$ con $N=80$.
    Cabe holgadamente con \textit{batch size} = 16 y FP16.
\end{itemize}

\paragraph{Arquitectura usada.}
Bloque Mamba con $d_{\text{model}} = 80$ y $N = 4$ capas apiladas,
seguido de \textit{average pooling} sobre el tiempo y una cabeza FC
lineal que produce el logit. Entrenado en FP16 con
\textit{gradient checkpointing} para ajustarse a la VRAM disponible.
Las hiperparametros exactos viven en \texttt{configs/mamba\_small.yaml}.

% ============================================================
\section{Protocolo experimental (8\,\%)}
\label{sec:protocol}
% RUBRIC EXCELENTE: protocolo sólido, test ciego, split por TIC, prevención de leakage.

\subsection{Splits por TIC ID}
% Justificar: una estrella en múltiples sectores -> mezclar = leakage severo.
% Proporción 70/15/15. Versionados en data/splits/.
% Conteos por clase ya están en Tabla \ref{tab:dataset-counts} de la Sec. 2.

\subsection{Test sellado}
% Política: una sola evaluación por modelo, post-tuning.
% Materializado en scripts/wsl2/eval_mamba_test_all.sh
% Cada modelo+seed consume UNA evaluación.

\subsection{Multi-seed como sustituto de K-fold}
% Justificación de costo: K=5 sobre Mamba en RTX 3050 4GB es prohibitivo
% (>= 25 h por modelo). Reportamos:
%   - Mamba single: 5 seeds (42, 123, 456, 789, 2024)
% Multi-seed da variabilidad estadística equivalente para el reporte.

\subsection{Prevención de leakage}
% (1) Splits por TIC ID (no por sector).
% (2) Normalización por curva individual (no estadísticas globales).
% (3) Augmentation solo en train.

% ============================================================
\section{Optimización de hiperparámetros (7\,\%)}
\label{sec:tuning}
% RUBRIC EXCELENTE: búsqueda sistemática + justificación costo + reproducibilidad.

\subsection{Sanity overfit por modelo}
% Cada modelo nuevo pasa un sanity sobre subset de 64 con dropout=0:
% debe llegar a val_auc=1.0 antes de entrenar real.
% Si no llega, se revisa preprocesamiento / arquitectura antes de gastar GPU.
\begin{table}[H]
\centering
\caption{Resultados de sanity overfit (subset N=64, sin dropout).}
\begin{tabular}{lcc}
\toprule
Modelo & Epoch hasta val\_auc=1.0 & Resultado \\
\midrule
CNN baseline                     & & \\
Mamba small                      & & \\
\bottomrule
\end{tabular}
\end{table}

\subsection{Sweep multi-seed}
\begin{table}[H]
\centering
\caption{Variabilidad por seed (val\_auc).}
\begin{tabular}{lc}
\toprule
Run & val\_auc \\
\midrule
mamba\_small\_locked        & 0.750 \\
mamba\_small\_seed42        & \\
mamba\_small\_seed123       & \\
mamba\_small\_seed456       & \\
mamba\_small\_seed789       & \\
mamba\_small\_seed2024      & \\
\midrule
\textbf{Mean $\pm$ std}     & \\
\bottomrule
\end{tabular}
\end{table}

\subsection{Justificación de no hacer grid/Optuna}
% Costo estimado: 1 epoch Mamba = ~5 min. 50 epochs x 20 configs = ~83 h en RTX 3050 4GB.
% Multi-seed sobre arquitectura locked es un trade-off documentado.
% La rúbrica acepta "búsqueda parcial pero válida"; la justificación de costo
% sube el ítem a "sistemática + justificación".

\subsection{Reproducibilidad de cada run}
% Cada run en experiments/ guarda: config.yaml, env_info.txt, git_info.txt,
% train.log, metrics.csv, checkpoints/{best,last}.pt.

% ============================================================
\section{Métricas y resultados (parte de 8\,\%)}
\label{sec:results}

\subsection{Métricas utilizadas}
% Justificar por qué NO accuracy (clase desbalanceada).
\begin{itemize}
  \item AUC-ROC: área bajo la curva ROC, invariante a threshold.
  \item AUC-PR: relevante para clase positiva minoritaria.
  \item F1, Recall, Precision @ threshold=0.5.
  \item Brier score: $\frac{1}{N}\sum(y_i - p_i)^2$, calibración.
\end{itemize}

% BCE balanceada (función de pérdida):
\begin{equation}
\mathcal{L}_{\text{BCE}} = -\frac{1}{N}\sum_{i=1}^{N}
  \bigl[ w_+ \, y_i \log p_i + (1-y_i) \log(1-p_i) \bigr],
\quad w_+ = \frac{N_{\text{FP}}}{N_{\text{CP}}}.
\end{equation}

\subsection{Tabla principal}
\begin{table}[H]
\centering
\caption{Métricas en test sellado para los cuatro modelos del estudio.}
\label{tab:main-results}
\begin{tabular}{lcccccc}
\toprule
Modelo & Input & AUC-ROC & AUC-PR & F1 & Recall & Precision \\
\midrule
Random estratificado        & ---           & & & & & \\
LogReg (catálogo)           & metadata      & & & & & \\
CNN single-branch           & global view   & & & & & \\
Mamba single (locked)       & global view   & & & & & \\
\textbf{Mamba ensemble 5s}  & global view   & \textbf{} & & & & \\
\bottomrule
\end{tabular}
\end{table}

\subsection{Curvas ROC comparativas}
\begin{figure}[H]
\centering
\safeincludegraphics[width=0.85\textwidth]{figures/roc_tier1.png}
\caption{Curvas ROC en test sellado. Mamba ensemble vs baselines.}
\label{fig:roc-tier1}
\end{figure}

\subsection{Matriz de confusión y calibración}
\begin{figure}[H]
\centering
\begin{subfigure}{0.48\textwidth}
  \safeincludegraphics[width=\textwidth]{results/mamba_ensemble/confusion_matrix.png}
  \caption{Matriz de confusión \@ threshold=0.5}
\end{subfigure}
\hfill
\begin{subfigure}{0.48\textwidth}
  \safeincludegraphics[width=\textwidth]{results/mamba_ensemble/calibration.png}
  \caption{Curva de calibración (reliability diagram)}
\end{subfigure}
\caption{Mamba ensemble: matriz de confusión y calibración.}
\end{figure}

% ============================================================
\section{Análisis de errores (parte de 8\,\%)}
\label{sec:errors}
% RUBRIC EXCELENTE: error analysis + variabilidad.

\subsection{Distribución de probabilidades por clase verdadera}
\begin{figure}[H]
\centering
\safeincludegraphics[width=0.7\textwidth]{results/error_analysis/mamba_ensemble/prob_histogram.png}
\caption{Histograma de \texttt{y\_prob} separado por clase verdadera. Idealmente bimodal.}
\end{figure}

\subsection{Tasa de error por feature física}
\begin{figure}[H]
\centering
\safeincludegraphics[width=0.95\textwidth]{results/error_analysis/mamba_ensemble/error_rate_by_feature.png}
\caption{Tasa de error vs bins de \texttt{pl\_orbper}, \texttt{pl\_trandep}, \texttt{st\_tmag}.
Detecta si el modelo falla más en períodos cortos, tránsitos profundos o estrellas brillantes.}
\end{figure}

\subsection{Top-10 falsos negativos y falsos positivos}
\begin{figure}[H]
\centering
\begin{subfigure}{0.48\textwidth}
  \safeincludegraphics[width=\textwidth]{results/error_analysis/mamba_ensemble/top_fn_curves.png}
  \caption{Top FN (CP clasificados como FP)}
\end{subfigure}
\hfill
\begin{subfigure}{0.48\textwidth}
  \safeincludegraphics[width=\textwidth]{results/error_analysis/mamba_ensemble/top_fp_curves.png}
  \caption{Top FP (FP clasificados como CP)}
\end{subfigure}
\caption{Casos más confidentes mal clasificados por Mamba ensemble.}
\end{figure}

% ============================================================
\section{Interpretabilidad --- XAI (7\,\%)}
\label{sec:xai}
% RUBRIC EXCELENTE: técnica adecuada + hallazgos claros y accionables.

\subsection{Métodos aplicados}
% Justificar: NO usamos "atención" porque Mamba no es un Transformer.
\begin{itemize}
  \item \textbf{Gradient saliency}: $s_i = \left|\dfrac{\partial \hat{y}}{\partial x_i}\right|$.
  \item \textbf{Integrated Gradients} (50 pasos, baseline $x'=\mathbf{0}$):
  \begin{equation}
  \text{IG}_i(x) = (x_i - x'_i)\int_0^1 \frac{\partial \hat{y}(x' + \alpha(x - x'))}{\partial x_i}\,d\alpha.
  \end{equation}
  Aproximado por suma de Riemann con 50 puntos.
  \item \textbf{Occlusion sensitivity}: $o_i = \hat{y}(x) - \hat{y}(x_{[i-w:i+w]\leftarrow 1.0})$.
\end{itemize}

\subsection{Política de selección de casos}
% 8 casos: top-2 más confidentes por cuadrante (TP, TN, FN, FP).
% Modelo analizado: Mamba best seed (789).
% Confianza = |y_prob - 0.5|.

\subsection{Hallazgos}
\begin{figure}[H]
\centering
\safeincludegraphics[width=0.95\textwidth]{figures/xai/mamba_seed789/_summary.png}
\caption{Summary XAI para Mamba seed 789 sobre 8 casos representativos.}
\label{fig:xai-summary}
\end{figure}

\begin{figure}[H]
\centering
\safeincludegraphics[width=0.85\textwidth]{figures/xai/mamba_seed789/TP_1_tic218795833_saliency.png}
\caption{Saliency sobre un TP confidente (TIC 218795833). El gradiente se concentra
en la región del dip de tránsito.}
\end{figure}

% Discusión accionable:
% (1) ¿Dónde mira Mamba? ¿En el dip o en el ruido del baseline?
% (2) Los FN ¿muestran patrón distinto (saliency disperso)? -> hipótesis para futuro work.

\begin{tcolorbox}[colback=blue!5,colframe=blue!50!black,title=Hallazgo XAI principal]
% TODO: 2-3 líneas con el insight principal del análisis.
\end{tcolorbox}

% ============================================================
\section{Ingeniería y reproducibilidad (5\,\%)}
\label{sec:repro}
% RUBRIC EXCELENTE: repo limpio + README completo + seeds + configs + ejecución robusta.

\subsection{Estructura del repositorio}
\begin{lstlisting}[caption={Árbol del repositorio (top-level).}]
mamba-exoplanet/
+-- configs/             # YAMLs por experimento (versionados)
+-- data/
|   +-- splits/          # CSVs train/val/test por TIC (versionados)
|   +-- raw/processed/   # gitignored
+-- src/exoplanet/       # paquete instalable (src layout)
|   +-- data/ models/ training/ evaluation/ utils/
+-- scripts/             # CLIs reproducibles (train, evaluate, preprocess)
|   +-- wsl2/            # helpers shell para entorno WSL2
+-- experiments/         # outputs por run (gitignored)
+-- paper/               # reporte + figuras + resultados
+-- notebooks/           # exploración numerada
+-- tests/               # 8 archivos, 49 tests
+-- docs/                # documentación interna
+-- pyproject.toml       # PEP 621 + ruff + pytest
+-- README.md            # 10 pasos para reproducir Etapa 2
\end{lstlisting}

\subsection{Configs YAML versionados}
% Mostrar fragmento de configs/mamba_small.yaml.

\subsection{Semillas y determinismo}
% set_seed() configura torch + numpy. Test test_seeds.py lo verifica.

\subsection{Artefactos por run}
\begin{itemize}
  \item \texttt{config.yaml} --- copia del YAML usado.
  \item \texttt{env\_info.txt} --- Python, torch, CUDA, mamba-ssm, OS.
  \item \texttt{git\_info.txt} --- commit hash + estado del working tree.
  \item \texttt{train.log} --- log epoch-by-epoch.
  \item \texttt{metrics.csv} --- métricas tabulares.
  \item \texttt{checkpoints/\{best,last\}.pt} --- pesos.
  \item \texttt{eval\_test/} --- métricas de test sellado (1 vez por modelo).
\end{itemize}

\subsection{Entornos: Windows nativo vs WSL2}
\begin{table}[H]
\centering
\caption{Entornos canónicos.}
\begin{tabular}{lll}
\toprule
Componente & Pipeline general (Windows) & Modelo Mamba (WSL2 Ubuntu 24.04) \\
\midrule
Python      & 3.11.9              & 3.11                          \\
torch       & 2.11.0+cu128        & 2.5.1+cu121                   \\
mamba-ssm   & ---                 & 2.2.6                          \\
CUDA driver & 581.83              & 12.1                           \\
GPU         & RTX 3050 4GB        & RTX 3050 4GB                   \\
\bottomrule
\end{tabular}
\end{table}

% ============================================================
\section{Suite de tests}
\label{sec:tests}
% Reusar docs/tests_overview.md.
% Total: 8 archivos, 49 tests.

\subsection{Estrategia general}
% (1) Smoke / contratos: test_smoke.py, test_seeds.py.
% (2) Componentes aislados: test_augment.py, test_collate.py,
%     test_metrics.py, test_cnn_baseline.py.
% (3) End-to-end con skip-si-no-hay-datos: test_dataset.py, test_training_smoke.py.

\begin{table}[H]
\centering
\caption{Cobertura de la suite de pruebas.}
\label{tab:tests}
\begin{tabular}{llc}
\toprule
Archivo & Cubre & Tests \\
\midrule
\texttt{test\_smoke.py}           & Imports del paquete y subpaquetes        & 2 \\
\texttt{test\_seeds.py}           & set\_seed() reproducible torch + numpy   & 3 \\
\texttt{test\_augment.py}         & 4 augmentations + Compose + registry     & 21 \\
\texttt{test\_collate.py}         & Batching con campos opcionales           & 4 \\
\texttt{test\_metrics.py}         & AUC/F1/Recall/Precision casos canónicos  & 4 \\
\texttt{test\_cnn\_baseline.py}   & Forward, backward, params, kwargs        & 4 \\
\texttt{test\_dataset.py}         & \texttt{LightCurveDataset}: schema y \textit{loading} & 5 \\
\texttt{test\_training\_smoke.py} & End-to-end con \texttt{configs/smoke.yaml}            & 3 \\
\addlinespace
Otros (\textit{loaders}, \textit{utils})  &                                  & 3 \\
\midrule
\textbf{Total}                    &                                          & \textbf{49} \\
\bottomrule
\end{tabular}
\end{table}

% Detallar bloque por bloque (ver docs/tests_overview.md).

% ============================================================
\section{Discusión}
\label{sec:discussion}
% (a) Mamba > CNN: confirma valor del SSM sobre conv pura.
% (b) Threshold del proyecto: AUC >= 0.88 era aspiracional; resultado real reportado honestamente.
% (c) Variabilidad multi-seed: ensemble estabiliza el reporte final.

% ============================================================
\section{Limitaciones y trabajo futuro}
\label{sec:limits}
% - Dataset 1576 ejemplos (chico vs 15K de Kepler).
% - Sin K-fold (justificado por costo de RTX 3050 4GB).
% - Sin transfer learning desde checkpoint Kepler.
% - Modelos multibranch (global + local) entrenados pero fuera del alcance de este reporte.
% - Agente LLM (vetting tool) -> Etapa 3.

% ============================================================
\section{Conclusiones}
\label{sec:conclusions}
% 1 párrafo. Qué se entrega, qué se aprendió, qué sigue.

% ============================================================
\bibliographystyle{plain}
\begin{thebibliography}{9}
\bibitem{gu2023mamba} Gu, A. \& Dao, T. (2023).
  \textit{Mamba: Linear-Time Sequence Modeling with Selective State Spaces}.
  arXiv:2312.00752.
\bibitem{shallue2018} Shallue, C. J. \& Vanderburg, A. (2018).
  \textit{Identifying Exoplanets with Deep Learning}. AJ 155(2), 94.
\bibitem{exominer} Valizadegan, H. et al. (2022). \textit{ExoMiner}. ApJ 926(2), 120.
\bibitem{lightkurve} Lightkurve Collaboration (2018). ascl:1812.013.
% TODO: DART-Vetter, ExoMiner++, NASA TESS archive, NASA Exoplanet Archive.
\end{thebibliography}

% ============================================================
\appendix
\section{Configs YAML completos}
\label{app:configs}
% Listing de configs/mamba_small.yaml.

\section{Tabla completa de runs}
\label{app:runs}
% Tabla con timestamp, modelo, seed, val_auc, test_auc, path en experiments/.

\end{document}
